\documentclass[10pt,
% a4paper,
%twocolumn,
fleqn,
%landscape, 
% papersize,
dvipdfmx,
uplatex
]{jsarticle}



\def\maru#1{\textcircled{\scriptsize#1}}%丸囲み番号

% \RequirePackage[2020/09/30]{platexrelease}

%太字設定
\usepackage[deluxe]{otf}

\usepackage{emathEy}

\usepackage[g]{esvect}

%大きな文字
\usepackage{fix-cm}

%定理環境
\usepackage{emathThm}
%\theoremstyle{boxed}
\theorembodyfont{\normalfont}
\newtheorem{Question}{問題}[subsection]
\newtheorem{Q}{}[subsection]
\newtheorem{question}[Question]{}
\newtheorem{quuestion}{}[subsection]

%セクション，大問番号のデザイン
\renewcommand{\labelenumi}{(\arabic{enumi})}
\renewcommand{\theenumii}{\alph{enumii})}
\renewcommand{\thesection}{第\arabic{section}章}

%用紙サイズの詳細設定
\usepackage{bxpapersize}
\papersizesetup{size={80mm,45mm}}
\usepackage[top=0.7zw,bottom=0truemm,left=3truemm,right=133truemm]{geometry}
\usepackage[dvipdfmx]{graphicx}

%余白など
\usepackage{setspace} % 行間
\setlength{\mathindent}{1zw}
\setlength\parindent{0pt}


%色カラーに関する設定
\usepackage{color}
\definecolor{shiro}{rgb}{0.95703125,0.87109375,0.7421875}
\definecolor{kin}{rgb}{0.95703125,0.87109375,0.7421875}
\definecolor{orange}{rgb}{1,0.7,0.2}
\definecolor{bradorange}{rgb}{1,0.5,0}
\definecolor{pink}{rgb}{0.9176,0.5686,0.5960}
\definecolor{mizu}{rgb}{0.6156,0.8,0.9955}
\color{kin}
% \pagecolor{hukamido}

\usepackage{at}%図の配置
% \usepackage{wallpaper}

\begin{document}



\at(0cm,0cm){\includegraphics[width=8cm,bb=0 0 1920 1080]{./youtube/thumbnails/templates/smart_background/２次関数.jpeg}}
{\color{orange}\bf\boldmath\LARGE\underline{区間に少なくとも$1$つの解}}\vspace{0.3zw}

\large 
\bf\boldmath 問．$x$についての$2$次方程式

\huge
\vspace{0.1zw}
\hspace{0.5zw}$x^2-2ax+a+2=0\vspace{0.5zw}$

\large 
の解が$1$$<$$x$$<$$3$の範囲に少なくとも$1$つ存在するような定数$a$の値の範囲を求めよ
\at(7.0cm,0.2cm){\small\color{bradorange}$\overset{\text{２次関数}}{\text{応用}}$}


\newpage



\at(0cm,0cm){\includegraphics[width=8cm,bb=0 0 1920 1080]{./youtube/thumbnails/templates/smart_background/２次関数.jpeg}}
{\color{orange}\bf\boldmath\Large\underline{区間で常に$2$次不等式が成立}}\vspace{0.3zw}

\large 
\bf\boldmath 問．$0\leqq x\leqq 3$を満たす$x$に対して，

\huge
\vspace{0.1zw}
\hspace{0.5zw}$x^2-2ax+a+2>0\vspace{0.5zw}$

\large 
が成り立つような\vspace{-0.zw}\\
\hfill 定数$a$の値の範囲を求めよ．
\at(7.0cm,0.2cm){\small\color{bradorange}$\overset{\text{２次関数}}{\text{応用}}$}


\newpage



\at(0cm,0cm){\includegraphics[width=8cm,bb=0 0 1920 1080]{./youtube/thumbnails/templates/smart_background/２次関数.jpeg}}
{\color{orange}\bf\boldmath\Large\underline{$「$すべて$」$と$「$ある$」$の交換}}\vspace{0.3zw}

\scriptsize 
\bf\boldmath 問．$x,\;y$についての条件$p$を次のように定める．

\vspace{0.6zw}
\hspace{0.5zw}$p:-x^2+\left(a-2\right)x+a-4<y<x^2-\left(a-4\right)x+3\vspace{0.3zw}$
\vspace{0.3zw}

\hfill 次の各々が成立するための$a$の値の範囲を求めよ．

\normalsize
(1)  どんな$x$に対しても，\vspace{-0.2zw}\\
\hfill 適当な$y$をとれば，$p$が成り立つ．\\
(2)  適当な$y$をとれば，\vspace{-0.2zw}\\
\hfill どんな$x$に対しても，$p$が成り立つ．\\

\at(7.0cm,0.2cm){\small\color{bradorange}$\overset{\text{２次関数}}{\text{応用}}$}


\newpage



\at(0cm,0cm){\includegraphics[width=8cm,bb=0 0 1920 1080]{./youtube/thumbnails/templates/smart_background/２次関数.jpeg}}
{\color{orange}\bf\boldmath\Large\underline{$2$変数の$「$すべて$」$と$「$ある$」$}}\vspace{0.3zw}

\scriptsize 
\bf\boldmath 問．$2$つの関数

\vspace{0.6zw}
\hspace{1zw}$f\left(x\right)=x^2+2x-2$,\hspace{2zw}$g\left(x\right)=-x^2+2x+a+1$\vspace{0.6zw}

\hfill について，次の各々が成立するような$a$の値の範囲を求めよ．

\normalsize
(1)  $-2\leqq x\leqq 2$を満たす\vspace{-0.2zw}\\
\hfill すべての$x_1,\;x_2$に対して$f\left(x_1\right)<g\left(x_2\right)$\\
(2)  $-2\leqq x\leqq 2$を満たす\vspace{-0.2zw}\\
\hfill ある$x_1,\;x_2$に対して$f\left(x_1\right)<g\left(x_2\right)$\\

\at(7.0cm,0.2cm){\small\color{bradorange}$\overset{\text{２次関数}}{\text{応用}}$}


\newpage



\at(0cm,0cm){\includegraphics[width=8cm,bb=0 0 1920 1080]{./youtube/thumbnails/templates/smart_background/２次関数.jpeg}}
{\color{orange}\bf\boldmath\LARGE\underline{ある実数$x$で成立}}\vspace{0.3zw}

\large 
\bf\boldmath 問．ある実数$x$に対して$2$次不等式

\Large
\vspace{0.3zw}
\hspace{0.5zw}$ax^2+\left(a-1\right)x+a-1<0\vspace{0.3zw}$


が成立するための\\
\hfill 定数$a$の範囲を求めよ．
\at(7.0cm,0.2cm){\small\color{bradorange}$\overset{\text{２次関数}}{\text{応用}}$}


\newpage



\at(0cm,0cm){\includegraphics[width=8cm,bb=0 0 1920 1080]{./youtube/thumbnails/templates/smart_background/２次関数.jpeg}}
{\color{orange}\bf\boldmath\LARGE\underline{すべての実数$x$で成立}}\vspace{0.3zw}

\large 
\bf\boldmath 問．すべての実数$x$に対して$2$次不等式


\Large
\vspace{0.3zw}
\hspace{0.3zw}$ax^2+\left(a-1\right)x+a-1<0\vspace{0.3zw}$


が成立するための\\
\hfill 定数$a$の範囲を求めよ．
\at(7.0cm,0.2cm){\small\color{bradorange}$\overset{\text{２次関数}}{\text{応用}}$}


\end{document}

